% ======================================================================
% Эксперимент 3.5.1. Численные производные и оптимальный шаг.
% Ответственный: A — Плешнев Арсений.
% Подраздел 3.5 «Численная проверка» начинается в report.tex; этот файл —
% первый эксперимент подраздела (заголовок subsection не нужен).
%
% Что сдать (main.tex: теория 1504–1516, пункт 1541–1544, ожидаемый результат
% 1563–1565, пример 1583–1624; требования к отчёту 1875–1879):
%   * функции численных частных производных и градиента (односторонние
%     и центральные разности) — в solution.py;
%   * график ошибки оценки dT/dlambda от шага h в логарифмическом масштабе;
%   * найденные по эксперименту оптимальные шаги для обеих схем;
%   * объяснение, почему ошибка не убывает неограниченно (остаток Тейлора
%     растёт с h, ошибка округления растёт как eps/h);
%   * указать точку, диапазон шагов, допуски; сравнить с оценками масштаба
%     оптимального шага из примера.
%
% Параметры варианта 3: рабочая точка из п. 3.1 (mu=80, lambda=600, c=8).
% Запуск: python experiments.py 3.5.1. Рисунок: figures/3-5-1-diff-error.pdf.
% ======================================================================

\taskhead{Эксперимент 3.5.1. Численные производные и оптимальный шаг}

\taskhead{Постановка эксперимента}
\todo{A: описать схемы (односторонняя и центральная разности), точку вычисления, диапазон шагов $h$, точное значение $T_\lambda$ для сравнения; указать допуски.}

\taskhead{Программный результат и сравнение}
\todo{A: листинг вызова своих функций из solution.py и запуск python experiments.py 3.5.1.}
\taskfigure[0.75]{figures/3-5-1-diff-error.pdf}{Ошибка оценки $\partial T/\partial\lambda$ от шага $h$ для обеих схем (логарифмические оси)}{fig:3-5-1}
\begin{table}[htbp]
\centering\small
\caption{Оптимальные шаги численного дифференцирования}
\label{tab:3-5-1}
\begin{tabular}{|l|r|r|}\hline
Схема & Оптимальный шаг $h$ & Минимальная ошибка\\\hline
Односторонняя разность & — & —\\
Центральная разность & — & —\\
\hline\end{tabular}
\end{table}

\taskhead{Вывод}
\todo{A: почему ошибка не убывает неограниченно при уменьшении $h$; сравнить найденные шаги с оценками масштаба $s\sqrt{\varepsilon}$ и $s\varepsilon^{1/3}$ (main.tex, строки 1616–1624).}
